% ECONOMICS_PROBLEM: {"id": "G12-290", "title": "Sources of asset-price volatility", "jel_code": "G12", "source_id": "OP-0311"}
% !TeX program = xelatex
\documentclass[11pt,a4paper]{article}
\usepackage[margin=23mm]{geometry}
\usepackage{fontspec}
\setmainfont{Latin Modern Roman}
\usepackage{amsmath,amssymb,mathtools}
\usepackage{xcolor,enumitem,booktabs,longtable,array,xurl,hyperref}
\definecolor{muted}{HTML}{53636C}
\hypersetup{colorlinks=true,urlcolor=blue,linkcolor=blue}
\setlength{\parindent}{0pt}
\setlength{\parskip}{5pt}
\newcommand{\R}{\mathbb R}
\newcommand{\E}{\mathbb E}
\newcommand{\N}{\mathbb N}
\newcommand{\ind}{\mathbf 1}
\newcommand{\Var}{\operatorname{Var}}
\newcommand{\Cov}{\operatorname{Cov}}
\newcommand{\supp}{\operatorname{supp}}
\DeclareMathOperator*{\argmin}{arg\,min}
\DeclareMathOperator*{\argmax}{arg\,max}
\newcommand{\entrymeta}[3]{{\sffamily\small\color{muted}\textbf{\texttt{#1}}\quad|\quad #2\par #3\par}}
\newcommand{\Problemref}[1]{\texttt{#1}}
\newcommand{\JELref}[2]{\texttt{#2} (JEL \texttt{#1})}
\begin{document}
\section*{Sources of asset-price volatility}
\textbf{Problem ID:} \texttt{G12-290}\quad \textbf{JEL:} \texttt{G12}\par
% BEGIN ECONOMICS STATEMENT
\label{problem:OP-0311}
{\sffamily\small\color{muted}Original subject group: Asset pricing and financial markets\par}
\label{definitions:problem:30007593}
\textbf{Audit classification:} Published research agenda. The constant-discount variance inequality is a proved conditional-expectation identity. The research agenda concerns richer mechanisms; variance-removal contrasts are editorial and generally nonadditive.
\subsection*{Definitions and precise research target}
Fix an asset paying nonnegative dividends \(D_t\), its ex-dividend price \(P_t\), an objective probability law \(\mathbb P\), and the information filtration \((\mathcal F_t)\). Assume the discounted dividend sum is square-integrable. Under the restricted constant-discount, rational-expectations benchmark, let
\[P_t^*=\sum_{k\geq1}\delta^kD_{t+k},\qquad P_t=E_{\mathbb P}[P_t^*\mid\mathcal F_t],\quad0<\delta<1.\]
For stationary, consistently detrended objects this implies \(\operatorname{Var}(P_t)\leq\operatorname{Var}(P_t^*)\); it does not apply automatically to changing discount rates. Enlarge the model to allow stochastic marginal utilities, heterogeneous subjective probability laws and dividend-growth news. The task is to identify the contribution of each mechanism to observed price innovations using dividend realizations, investor-expectation data and independently measured financing conditions. Define the contribution as a counterfactual innovation variance when one mechanism is fixed at its conditional baseline and all remaining equilibrium conditions are solved again. Specify the baseline, detrending and terminal-value assumptions explicitly. Determine whether the joint observations identify these counterfactual shares, or only an identified set. Correlations between components must be reported rather than silently assigned to one channel. A rejection of the restricted benchmark is not a proof of irrationality or of an unsolved variance inequality.

\textbf{Additional formal definitions and scope (editorial).} The sum defining \(P_t^*\) converges in \(L^2(\mathbb P)\), and conditional expectation is with respect to the same law and information as pricing. The variance bound holds at each fixed date by the law of total variance; stationarity is needed for pooling dates, not for that identity. A trend-removal operation must preserve the stated conditional-expectation relation before applying the bound to transformed data. In a richer structural model \(\theta\), define innovation \(\nu_t=P_t-E[P_t\mid\mathcal F_{t-1}]\). For channel \(k\), let \(\nu_t^{(-k)}\) be the innovation under an explicit intervention replacing its equation by a baseline equation, with other equations and shock coupling stated. Define \(V_k=\operatorname{Var}(\nu_t)-\operatorname{Var}(\nu_t^{(-k)})\); \(V_k/\operatorname{Var}(\nu_t)\) is a variance difference, potentially negative, and need not sum to one across channels. An additive decomposition instead requires an identified orthogonal-shock representation or a declared allocation of interaction terms.

For the identification part, declare an admissible class \(\Theta\) of structural models, including preferences, technologies, shock laws, measurement/sampling rules and any equilibrium selector. If \(O\) denotes the complete observable history and \(T(\theta)\) the target defined above, the identified set is \(\mathcal I_T=\{T(\theta):\theta\in\Theta,\ \mathcal L_\theta(O)=\mathcal L_{\rm obs}(O)\}\); point identification means this set is a singleton. A \(do(a)\) comparison replaces stated policy/technology equations while fixing the declared exogenous law and re-solving the remaining equations. These definitions do not assert that an identification theorem holds without further restrictions.
\subsection*{Variants and assumption changes}
\begin{enumerate}
\item \textbf{Constant discounting (established benchmark).} Under \(P=E[P^*\mid\mathcal F]\) and \(P^*\in L^2\), prove/test the fixed-date variance restriction; the inequality itself is resolved.
\item \textbf{Stochastic discounting and beliefs (editorial).} Specify independent measures of cash-flow forecasts and discount rates; identify variance differences or an orthogonal news decomposition, allowing negative contrasts and interactions.
\end{enumerate}
\subsection*{References and source audit}
\begin{enumerate}
\item Primary AEA journal scan verifies title, author and 1981 publication; no DOI supplied or invented. Locator: Opening, p.421; variance discussion, p.423; Section VI, pp.433--434. Support: Restricted variance benchmark. Full scan accessible. URL: \url{https://www.aeaweb.org/aer/top20/71.3.421-436.pdf}.
\item Author-hosted published chapter and NBER/publisher record verify authors, 2024 publication and DOI; NBER Macroeconomics Annual 2023 is the volume label. Locator: Section VI, printed pp.341--342 (PDF pp.32--33). Support: Expectations research agenda. Full PDF accessible; DOI separately verified at NBER/publisher. URL: \url{https://shleifer.scholars.harvard.edu/sites/g/files/omnuum10626/files/shleifer/files/long-term_expectations_and_aggregate_fluctuations.pdf}.
\end{enumerate}
\begin{enumerate}\raggedright
\item\hypertarget{definitions:source:30007593:1}{} Robert J. Shiller. 1981. \emph{Do Stock Prices Move Too Much to Be Justified by Subsequent Changes in Dividends?}. Opening, p.421; variance discussion, p.423; Section VI, pp.433--434.
\url{https://www.aeaweb.org/aer/top20/71.3.421-436.pdf}.
\item\hypertarget{definitions:source:30007593:2}{} Pedro Bordalo; Nicola Gennaioli; Rafael La Porta; Matthew O'Brien; Andrei Shleifer. 2024. \emph{Long-Term Expectations and Aggregate Fluctuations}. Section VI, printed pp.341--342 (PDF pp.32--33).
\url{https://shleifer.scholars.harvard.edu/sites/g/files/omnuum10626/files/shleifer/files/long-term_expectations_and_aggregate_fluctuations.pdf}.
DOI: \url{https://doi.org/10.1086/729198}.
\end{enumerate}

\subsection*{Common conventions: Source mathematical conventions}

These conventions apply to each entry unless explicitly replaced.

\textbf{Models and identification.} A primitive model \(m\) specifies all probability laws, feasible choices, information, equilibrium and measurement rules used by its target. The admissible class \(\mathcal M\) is fixed independently of outcomes. If \(P_O(m)\) is its observable probability law and \(\tau(m)\) the target, its identified set at observable law \(P\) is
\[
 \mathcal I_\tau(P)=\{\tau(m):m\in\mathcal M,\ P_O(m)=P\}.
\]
Point identification means that this set is a singleton. An empty set signals incompatibility of data and assumptions. Coordinate bounds need not describe the joint set. Bounds are sharp only if every retained target is attainable or arbitrarily approachable by compatible models. Finite bounds require explicit restrictions on unobserved quantities; unrestricted confounding, hidden wealth, extrapolation or arbitrary equilibrium selection can yield uninformative sets.

\textbf{Causal interventions.} A potential outcome \(Y_i(p)\) is the outcome under a complete policy regime \(p\), including timing, financing, information and market spillovers. The target population and units are fixed before treatment. With interference, use the complete assignment vector or a justified exposure mapping; individual no-interference notation alone cannot identify an economy-wide intervention. A difference of mean potential outcomes does not require the cross-world joint distribution. Conditioning on post-treatment migration, survival, enrollment or employment changes the estimand and needs separate justification.

\textbf{Statistical statements.} An estimator, confidence set, forecast or test specifies a sampling law, model class and loss/coverage criterion. Uniform \((1-\alpha)\)-coverage means \(\inf_{m\in\mathcal M}\Pr_m\{\tau(m)\in C_n\}\geq1-\alpha\), or an explicitly asymptotic version. The whole parameter space trivially has coverage, so informativeness requires a stated width, risk or power objective. A forecasting improvement is relative to a named benchmark, information set and evaluation population.

\textbf{Equilibrium and optimization.} An equilibrium comprises feasible plans, optimality, beliefs where needed, budget constraints and all market/resource clearing. The primitives or a stated benchmark define each correspondence \(\mathcal E(p;m)\). Nonemptiness is assumed or proved, never inferred just from compact policy sets. A selected-equilibrium target uses a declared selection rule; a robust target quantifies over all permitted equilibria. Use suprema or infima unless attainment is established. An \(\varepsilon\)-optimal policy is feasible and within \(\varepsilon\) of the finite optimum value. Report an empty feasible set separately.

\textbf{Welfare and accounting.} Normative utility, interpersonal normalization, social weights, numeraire, discounting and the use of public receipts are primitives. Behavioral utility need not coincide with the welfare criterion. Household welfare includes ownership income and rents once; adding profits again to compensating variation generally double-counts them. A monetary equivalent is defined by an explicit transfer and price convention, with monotonicity and range conditions. Average effects, mechanism contributions, transfers and welfare losses are different targets. Infinite sums require convergence and dynamic feasible plans require terminal or no-Ponzi conditions.

\textbf{Source versions.} A working-paper date, revision date and journal publication year are distinguished. A DOI for a journal version is not automatically the DOI of an earlier manuscript. Locators refer to the stated version and printed pages where specified. A metadata-only check does not verify a claim about a full-text conclusion. Every entry's source subsection narrows the provenance claim accordingly.
% END ECONOMICS STATEMENT
\end{document}
